\begin{table}[htbp]
\centering
\small
\setlength{\tabcolsep}{4pt}
\caption{Prueba de Hausman: efectos fijos frente a efectos aleatorios}
\label{tab:tabla5_hausman}
\begin{tabular}{l>{\centering\arraybackslash}p{2.6cm}>{\centering\arraybackslash}p{2.6cm}>{\centering\arraybackslash}p{2.6cm}}
\toprule
Coeficiente & Efectos fijos & Efectos aleatorios & Diferencia \\
\midrule
Escudo fiscal / activos (\%) & 47,944 & 48,636 & $-$0,692 \\
Rentabilidad del activo, ROA (\%) & $-$0,170 & $-$0,174 & 0,004 \\
Riesgo país, EMBIG (\%) & 2,228 & 2,274 & $-$0,046 \\
Tasa legal del IR (\%) & $-$1,691 & $-$1,659 & $-$0,032 \\
Hausman clásica, chi² (4) & 16,57 & p = 0,002 &  \\
Mundlak robusta, chi² (2) & 3,31 & p = 0,191 &  \\
\bottomrule
\end{tabular}
\par\vspace{2pt}\footnotesize\raggedright \textit{Nota.} Variable dependiente: D/V (\%); 1015 observaciones de 20 emisoras. Efectos aleatorios por MCG factibles (Swamy-Arora). La prueba clásica compara los cuatro coeficientes con varianzas convencionales; la versión de Mundlak contrasta las medias de empresa del escudo y del ROA con errores agrupados por empresa. H0: los efectos de empresa no están correlacionados con los regresores (efectos aleatorios consistentes).
\end{table}
